\chapter{Diseases of the Muscles}
\label{ch:muscles}
Skeletal muscles constitute approximately 40\% of body weight and are responsible for movement, posture, heat generation, and protection of internal organs. Muscle diseases (myopathies) encompass a heterogeneous group of disorders affecting skeletal muscle, characterized by progressive weakness, myalgia, and elevated serum creatine kinase (CK). These conditions may be genetic, autoimmune, infectious, metabolic, or toxic in etiology.

%-------------------------------------------------------------
\section{Channelopathies and Metabolic Myopathies}
%-------------------------------------------------------------

%-------------------------------------------------------------

%-------------------------------------------------------------

%-------------------------------------------------------------

%-------------------------------------------------------------

%-------------------------------------------------------------

%-------------------------------------------------------------

%-------------------------------------------------------------

\label{sec:Channelopathies and Metabolic Myopathies}
%-------------------------------------------------------------%-------------------------------------------------------------

\subsection{McArdle Disease}
\label{sec:McArdle Disease}
McArdle disease (glycogen storage disease type V) is an autosomal recessive metabolic myopathy caused by deficiency of muscle glycogen phosphorylase, the enzyme that breaks down glycogen to glucose-1-phosphate in skeletal muscle. Clinical features include exercise intolerance, muscle pain, and cramping during vigorous activity, often with a "second wind" phenomenon (improvement in symptoms after a brief rest during exercise); rhabdomyolysis may occur with severe exercise. The condition results from PYGM gene mutations leading to inability to mobilize glycogen stores during anaerobic metabolism. Diagnosis is based on elevated serum CK, ischemic forearm exercise test (failure of lactate to rise), and genetic testing (PYGM mutations); muscle biopsy shows glycogen accumulation and absent phosphorylase activity. Management includes exercise conditioning, high-protein diet, and avoidance of strenuous exercise; oral sucrose before exercise may improve exercise capacity. Prognosis is variable but many patients maintain functional independence with lifestyle modifications.

\subsection{Myasthenia Gravis}
\label{sec:Myasthenia Gravis}
Myasthenia gravis (MG) is a neuromuscular junction disorder caused by autoantibodies against components of the postsynaptic acetylcholine receptor (AChR) or, less commonly, muscle-specific kinase (MuSK) or lipoprotein-related protein 4 (LRP4). Prevalence is approximately 20 per 100,000. The condition results from impaired neuromuscular transmission due to antibody-mediated blockade or destruction of AChR. Clinical features include fluctuating, fatigable weakness that worsens with activity and improves with rest; ocular muscles are most commonly affected first (ptosis, diplopia), followed by bulbar muscles (dysphagia, dysarthria, facial weakness), and limb and respiratory muscles. Myasthenic crisis (acute respiratory failure) occurs in 15--20\% of patients. Diagnosis is based on AChR or MuSK antibody testing, edrophonium (Tensilon) test (transient improvement), repetitive nerve stimulation (decremental response), and single-fiber EMG (increased jitter). Thymoma is present in 10--15\% of patients. Management includes acetylcholinesterase inhibitors (pyridostigmine), immunosuppression (prednisone, azathioprine, mycophenolate), thymectomy, and targeted therapies (eculizumab, ravulizumab [complement inhibitors], efgartigimod, rozanolixizumab [FcRn inhibitors]). Prognosis is generally good with treatment, though myasthenic crisis remains a life-threatening complication.

\subsection{Periodic Paralysis}
\label{sec:Periodic Paralysis}
Periodic paralysis is a rare autosomal dominant channelopathy characterized by episodic attacks of muscle weakness or paralysis. Hypokalemic periodic paralysis presents with attacks of flaccid paralysis with low serum potassium, triggered by carbohydrate-rich meals, rest after exercise, or stress. Hyperkalemic periodic paralysis features attacks with elevated potassium, triggered by fasting or exercise. Andersen-Tawil syndrome combines periodic paralysis with cardiac arrhythmias and skeletal developmental abnormalities. The condition results from mutations in ion channel genes (CACNA1S, SCN4A, KCNJ2). Clinical features include attacks lasting hours to days; between attacks, patients may be asymptomatic or have mild fixed weakness. Diagnosis is based on genetic testing, potassium levels during attacks, and provocation testing. Management includes potassium supplementation during attacks and acetazolamide or dichlorphenamide prophylaxis for hypokalemic type; inhaled beta-agonists, glucose, and avoidance of fasting for hyperkalemic type. Prognosis is variable; attacks are typically reversible with appropriate management.

%-------------------------------------------------------------
\section{Inflammatory Myopathies}
%-------------------------------------------------------------

%-------------------------------------------------------------

%-------------------------------------------------------------

%-------------------------------------------------------------

%-------------------------------------------------------------

%-------------------------------------------------------------

%-------------------------------------------------------------

%-------------------------------------------------------------

\label{sec:Inflammatory Myopathies}
%-------------------------------------------------------------%-------------------------------------------------------------

\subsection{Dermatomyositis}
\label{sec:Dermatomyositis}
Dermatomyositis is an idiopathic inflammatory myopathy with characteristic skin manifestations, discussed in Chapter~\ref{ch:skin}. In addition to proximal muscle weakness, it presents with heliotrope rash, Gottron's papules, V-sign, shawl sign, and mechanic's hands. It is mediated by complement attack on intramuscular blood vessels (perifascicular atrophy on biopsy). Adults with dermatomyositis have an increased risk of underlying malignancy, necessitating age-appropriate cancer screening.

\subsection{Inclusion Body Myositis}
\label{sec:Inclusion Body Myositis}
Inclusion body myositis (IBM) is the most common acquired myopathy in adults over 50 and the most common inflammatory myopathy in men. It is characterized by slowly progressive, asymmetric muscle weakness affecting both proximal and distal muscles, with characteristic involvement of wrist and finger flexors (difficulty gripping objects) and quadriceps (difficulty rising from chairs and frequent falls); dysphagia occurs in 30--40\%. The condition results from a combination of inflammatory and degenerative mechanisms. Diagnosis is based on muscle biopsy showing endomysial inflammation, rimmed vacuoles, and intracellular amyloid deposits; CK is only mildly elevated. Management is largely supportive, with physical therapy to prevent falls and maintain function; the disease is refractory to immunosuppressive therapy. Prognosis is slowly progressive with significant functional impairment over time.

\subsection{Polymyositis}
\label{sec:Polymyositis}
Polymyositis is an idiopathic inflammatory myopathy characterized by symmetric proximal muscle weakness and endomysial CD8+ T-cell infiltration of muscle fibers. It primarily affects adults, with a slight female predominance. The condition results from autoimmune-mediated muscle fiber destruction. Clinical features include progressive symmetric proximal muscle weakness (difficulty rising from chairs, climbing stairs, lifting objects overhead), dysphagia, and rarely respiratory muscle weakness. Diagnosis is based on elevated CK (10--50 times normal), EMG showing myopathic changes, MRI showing muscle edema, and definitive muscle biopsy. Muscle-specific autoantibodies (anti-Jo-1, anti-Mi-2, anti-MDA5) may support diagnosis. Management includes corticosteroids and steroid-sparing immunosuppressive agents (methotrexate, azathioprine, mycophenolate); IVIG may be used for refractory disease. Prognosis is variable, with many patients responding to immunosuppression.

%-------------------------------------------------------------
\section{Muscle Injuries and Conditions}
%-------------------------------------------------------------

%-------------------------------------------------------------

%-------------------------------------------------------------

%-------------------------------------------------------------

%-------------------------------------------------------------

%-------------------------------------------------------------

%-------------------------------------------------------------

%-------------------------------------------------------------

\label{sec:Muscle Injuries and Conditions}
%-------------------------------------------------------------%-------------------------------------------------------------

\subsection{Compartment Syndrome}
\label{sec:Compartment Syndrome}
Compartment syndrome is elevated pressure within a closed fascial compartment, compromising tissue perfusion and leading to ischemia and necrosis of muscle and nerves. Acute compartment syndrome is a surgical emergency most commonly caused by fractures (tibial fractures most common), crush injuries, reperfusion injury, or tight casts. Clinical features include pain out of proportion to the injury, worsened by passive stretching of the affected muscles, with tense, swollen compartments; the five P's (pulselessness, paresthesia, paralysis, pallor, poikilothermia) are late signs. Diagnosis is based on clinical examination and compartment pressure measurement (pressures >30 mmHg or within 30 mmHg of diastolic pressure indicate need for fasciotomy). Management is emergent fasciotomy to decompress all affected compartments; delayed treatment leads to irreversible muscle necrosis, contracture (Volkmann's ischemic contracture), and potentially limb loss. Prognosis is excellent with early fasciotomy but poor if treatment is delayed.

\subsection{Fibromyalgia}
\label{sec:Fibromyalgia}
Fibromyalgia is a chronic pain condition characterized by widespread musculoskeletal pain, fatigue, sleep disturbance, cognitive dysfunction ("fibro fog"), and mood disorders, affecting 2--4\% of the population, predominantly women. The condition results from central sensitization (amplified pain processing in the CNS), altered neurotransmitter levels (substance P, serotonin), and abnormal functional connectivity in brain pain networks. Clinical features include chronic widespread pain, tender points (18 defined anatomic sites historically), fatigue, and sleep disturbance. Diagnosis is based on the 2010 ACR criteria using a widespread pain index and symptom severity scale; it is a clinical diagnosis of exclusion with no specific laboratory tests. Management is multimodal: exercise (most evidence-based intervention), cognitive-behavioral therapy, pharmacotherapy (pregabalin, duloxetine, milnacipran, amitriptyline, tramadol), and patient education. Prognosis is variable; symptoms are chronic but functional status can improve with multimodal management.

\subsection{Rhabdomyolysis}
\label{sec:Rhabdomyolysis}
Rhabdomyolysis is a syndrome of acute skeletal muscle necrosis with release of intracellular muscle contents (myoglobin, CK, potassium, phosphate, uric acid) into the circulation. Causes include crush injuries, prolonged immobilization, extreme exertion, seizures, malignant hyperthermia, medications (statins, fibrates, antipsychotics), electrolyte imbalances, infections, and inflammatory myopathies. Clinical features include myalgia, muscle weakness, dark (cola-colored) urine, and edema of affected muscles; complications include acute kidney injury (myoglobin nephrotoxicity), hyperkalemia (cardiac arrhythmias), disseminated intravascular coagulation, and compartment syndrome. Diagnosis is based on elevated CK (typically >5 times the upper limit of normal, often >10,000 IU/L), elevated myoglobin, and elevated transaminases. Management involves aggressive IV fluid resuscitation (target urine output 200--300 mL/h), urinary alkalinization (to prevent myoglobin precipitation), correction of electrolyte abnormalities, and monitoring for complications; renal replacement therapy may be necessary for severe AKI. Prognosis depends on severity and promptness of treatment; most patients recover with adequate hydration.

%-------------------------------------------------------------
\section{Muscular Dystrophies}
%-------------------------------------------------------------

%-------------------------------------------------------------

%-------------------------------------------------------------

%-------------------------------------------------------------

%-------------------------------------------------------------

%-------------------------------------------------------------

%-------------------------------------------------------------

%-------------------------------------------------------------

\label{sec:Muscular Dystrophies}
%-------------------------------------------------------------%-------------------------------------------------------------

\subsection{Becker Muscular Dystrophy}
\label{sec:Becker Muscular Dystrophy}
Becker muscular dystrophy (BMD) is an X-linked recessive disorder caused by in-frame mutations in the dystrophin gene, producing a truncated but partially functional dystrophin protein. It affects approximately 1 in 18,000--30,000 males. The condition results from dystrophin deficiency that is less severe than DMD. Clinical features include proximal muscle weakness with later onset (late childhood to early adulthood) and slower progression compared to DMD; patients may remain ambulatory into their 30s or beyond. Cardiomyopathy is a significant cause of morbidity and mortality, often presenting before significant skeletal muscle weakness. Diagnosis is based on genetic testing and muscle biopsy. Management is similar to DMD, with emphasis on cardiac surveillance and management. Prognosis is better than DMD, with longer ambulation and survival.

\subsection{Duchenne Muscular Dystrophy}
\label{sec:Duchenne Muscular Dystrophy}
Duchenne muscular dystrophy (DMD) is the most common and severe childhood muscular dystrophy, affecting approximately 1 in 3500--5000 live-born males. It is an X-linked recessive disorder caused by mutations in the dystrophin gene (DMD) on Xp21, leading to absence of the dystrophin protein essential for sarcolemmal integrity. Clinical features include progressive proximal muscle weakness beginning between ages 2--5, with Gowers' sign (using hands to "climb up" the thighs when rising from the floor), waddling gait, toe walking, and calf pseudohypertrophy; by age 10--12, most patients lose ambulation. Cardiomyopathy and respiratory failure develop in adolescence. Diagnosis is based on elevated CK (10--100 times normal), genetic testing showing dystrophin mutations, and muscle biopsy showing dystrophic changes with absent dystrophin on immunostaining. Management includes corticosteroids (deflazacort, prednisone) to slow progression, cardiac management (ACE inhibitors, beta-blockers), respiratory support, and newer therapies including exon-skipping agents (eteplirsen, golodirsen, viltolarsen, casimersen for specific mutations) and gene therapy (delandistrogene moxeparvovec). Prognosis is severe, with most patients losing ambulation by age 10--12 and survival into the third decade with respiratory support.

\subsection{Facioscapulohumeral Muscular Dystrophy}
\label{sec:Facioscapulohumeral Muscular Dystrophy}
Facioscapulohumeral muscular dystrophy (FSHD) is an autosomal dominant myopathy caused by contraction of D4Z4 repeats on chromosome 4q35, leading to inappropriate expression of the DUX4 gene. It has a prevalence of approximately 1 in 8000--20,000. Clinical features include initial involvement of facial muscles (inability to fully close the eyes, difficulty whistling, transverse smile) and scapular stabilizers ("scapular winging"), followed by humeral and pelvic girdle muscles; asymmetric involvement is common. Hearing loss and retinal telangiectasias (Coats disease) may occur. Diagnosis is based on genetic testing. Management includes physical therapy and orthopedic interventions (scapular bracing/fusion) and respiratory monitoring; there is no curative treatment. Prognosis is variable, with slowly progressive weakness but usually normal life expectancy.

\subsection{Lambert-Eaton Myasthenic Syndrome (LEMS)}
\label{sec:Lambert-Eaton Myasthenic Syndrome (LEMS)}
Lambert-Eaton myasthenic syndrome (LEMS) is a rare autoimmune disorder of the neuromuscular junction characterized by presynaptic impairment of acetylcholine release. Driven by pathogenic autoantibodies directed against voltage-gated P/Q-type calcium channels (VGCC) on presynaptic motor nerve terminals, the condition reduces calcium influx during nerve action potentials, preventing acetylcholine vesicle exocytosis into the synaptic cleft. Strongly associated with underlying malignancy (paraneoplastic LEMS, 50--60\% of cases, predominantly small cell lung cancer) or primary autoimmune disease. Clinical features present as proximal muscle weakness (affecting lower limbs first, causing difficulty rising from a chair), hyporeflexia or areflexia, and prominent autonomic dysfunction (dry mouth, erectile dysfunction, constipation). A hallmark feature is transient post-exercise facilitation (improvement of strength and reflexes after brief muscle contraction). Diagnosis is confirmed by high-titer serum anti-VGCC antibodies and repetitive nerve stimulation (RNS) demonstrating low baseline compound muscle action potential (CMAP) amplitude with dramatic increment ($>100\%$) following high-frequency stimulation or exercise. Management involves screening for small cell lung cancer, 3,4-diaminopyridine (amifampridine, enhances ACh release), pyridostigmine, and immunosuppression (prednisone, IVIG).

\subsection{Malignant Hyperthermia}
\label{sec:Malignant Hyperthermia}
Malignant hyperthermia (MH) is a life-threatening pharmacogenetic skeletal muscle disorder triggered by exposure to volatile halogenated anesthetic agents (isoflurane, sevoflurane, desflurane) or depolarizing muscle relaxants (succinylcholine) in susceptible individuals. Driven by autosomal dominant gain-of-function mutations in the RYR1 gene (encoding the ryanodine receptor 1 calcium release channel of the sarcoplasmic reticulum) or CACNA1S gene, volatile agents trigger uncontrolled, massive calcium efflux from the sarcoplasmic reticulum into the myoplasm. Persistent intracellular calcium accumulation drives sustained, uncoordinated muscle contractions, hypermetabolism, ATP depletion, and rhabdomyolysis. Clinical features present as rapid masseter muscle spasm (lockjaw), hypercapnia (rising end-tidal CO2, earliest and most sensitive sign), sinus tachycardia, muscle rigidity, severe mixed respiratory/metabolic acidosis, rhabdomyolysis (myoglobinuria), hyperkalemia, and late fulminant hyperthermia ($>42^\circ\text{C}$). Diagnosis is clinical during anesthesia; susceptibility is confirmed by caffeine-halothane contracture test (CHCT) or genetic testing. Management requires immediate discontinuation of triggering agents, hyperventilation with $100\%$ oxygen, administration of IV dantrolene sodium (ryanodine receptor antagonist), active cooling, and aggressive treatment of hyperkalemia and acidosis.

\subsection{Myotonic Dystrophy}
\label{sec:Myotonic Dystrophy}
Myotonic dystrophy (DM) is the most common muscular dystrophy in adults, with an autosomal dominant inheritance pattern. Type 1 (DM1, Steinert disease) results from a CTG trinucleotide repeat expansion in the DMPK gene on chromosome 19; Type 2 (DM2, proximal myotonic myopathy) results from a CCTG expansion in the CNBP gene. Clinical features in DM1 include distal and facial muscle weakness, myotonia (delayed muscle relaxation after contraction), cataracts (posterior subcapsular "Christmas tree" cataracts), cardiac conduction abnormalities, endocrine dysfunction, and cognitive impairment. Genetic anticipation (progressively earlier onset and more severe disease in successive generations) occurs. Diagnosis is based on genetic testing. Management is symptomatic: mexiletine for myotonia, cardiac pacemaker/defibrillator for conduction defects, and management of endocrine and other systemic manifestations. Prognosis is variable, with progressive disability and reduced life expectancy due to cardiac and respiratory complications.

